\section*{Chapitre \ref{ch:b2:metric} --- Topologie des espaces métriques}

\begin{solution}{exo:b2:metric:1}
$\delta$~: la symétrie et la séparation sont claires~; inégalité
triangulaire~: $\min(1, u + v) \leq \min(1,u) + \min(1,v)$ pour $u, v
\geq 0$ (si l'un des min vaut $1$, le membre de droite est $\geq 1$~;
sinon il vaut $u + v$). $d$~: c'est le tiré en arrière de
$\abs{\cdot}$ par l'injective $\arctan$~: les trois axiomes se
transfèrent.

Convergence~: pour $\delta$, $\delta(x_n, x) \to 0 \iff \abs{x_n - x}
\to 0$ (pour les petites valeurs les deux distances coïncident)~:
mêmes suites convergentes que d'habitude. Pour $d$~: $d(x_n, x) \to 0
\iff \arctan x_n \to \arctan x \iff x_n \to x$ (continuité et stricte
monotonie de $\arctan$ et de sa réciproque sur les plages
concernées)~: à nouveau la convergence usuelle.

$(\R, d)$ n'est \emph{pas} \omterm{def:b2:metric:complete}{complet}~: $x_n = n$ vérifie $d(x_p, x_q)
= \abs{\arctan p - \arctan q} \to 0$ (les deux tendent vers
$\frac\pi2$)~: de Cauchy~; mais $(x_n)$ ne converge pas pour $d$ (sa
$d$-limite serait une limite ordinaire). La complétude est une
propriété de la \emph{distance}, non seulement des suites
convergentes.
\end{solution}

\begin{solution}{exo:b2:metric:2}
Convergente $\Rightarrow$ de Cauchy~: $d(x_p, x_q) \leq d(x_p, \ell)
+ d(\ell, x_q)$. Bornée~: au-delà de $N$, $d(x_n, \ell) \leq 1$~; les
quelques premiers termes sont eux aussi dans un certain rayon.

De Cauchy $+$ sous-suite convergente $x_{\varphi(n)} \to \ell$~:
étant donné $\varepsilon$, pour $n$ grand, $d(x_n, \ell) \leq d(x_n,
x_{\varphi(n)}) + d(x_{\varphi(n)}, \ell) \leq 2\varepsilon$ (le
premier terme par Cauchy, puisque $\varphi(n) \geq n$).

\omterm{def:b2:metric:compact}{Compact} $\Rightarrow$ \omterm{def:b2:metric:complete}{complet}~: une suite de Cauchy a une sous-suite
convergente (compacité), donc converge.
\end{solution}

\begin{solution}{exo:b2:metric:3}
$d_\infty(f, g) = \sup_{\intcc{0}{1}} \abs{x - x^2} = \frac14$
(maximum de $x - x^2$ en $x = \frac12$).

$\overline B(0, 1) = \{f : \sup\abs f \leq 1\}$~: les fonctions
\omterm{def:b2:metric:continuity}{continues} à valeurs dans $\intcc{-1}{1}$.

$\{f : f(0) = 0\}$ est l'image réciproque de $\{0\}$ par
l'\emph{évaluation} $f \mapsto f(0)$, qui est $1$-lipschitzienne
($\abs{f(0) - g(0)} \leq d_\infty(f,g)$), donc \omterm{def:b2:metric:continuity}{continue}~: l'ensemble
est fermé (\cref{thm:b2:metric:globalcontinuity}). De même $\{f :
f(0) > 0\}$ est l'image réciproque de l'\omterm{def:b2:metric:topology}{ouvert}
$\intoo{0}{+\infty}$~: \omterm{def:b2:metric:topology}{ouvert}.
\end{solution}

\begin{solution}{exo:b2:metric:4}
\emph{\omterm{def:b2:metric:complete}{Complet}~:} une suite de Cauchy avec $\varepsilon = \frac12$ est
constante à partir d'un certain rang, donc convergente.

\emph{\omterm{def:b2:metric:compact}{Compact} ssi fini~:} si $X$ est fini, toute suite prend une
certaine valeur une infinité de fois (sous-suite constante). Si $X$
est infini, une suite de points deux à deux distincts a toutes ses
distances mutuelles égales à $1$~: aucune sous-suite de Cauchy, donc
aucune convergente.

\emph{Parties \omterm{def:b2:metric:connected}{connexes}~:} les singletons (et $\emptyset$). Tout $A$
ayant deux points $x \neq y$ se scinde en $\{x\} \cup (A
\setminus\{x\})$, tous deux \omterm{def:b2:metric:topology}{ouverts} dans $A$ (toute partie d'un
espace discret est ouverte --- les boules de rayon $\frac12$ sont des
singletons).
\end{solution}

\begin{solution}{exo:b2:metric:5}
La fonction $g(x) = d(x, f(x))$ est \omterm{def:b2:metric:continuity}{continue} sur le \omterm{def:b2:metric:compact}{compact} $X$
($\abs{g(x) - g(y)} \leq 2d(x,y)$ par deux inégalités
triangulaires), donc elle atteint son minimum en un certain $a$
(\cref{thm:b2:metric:compactprops}). Si $f(a) \neq a$~:
\[
g\bigl(f(a)\bigr) = d\bigl(f(a), f(f(a))\bigr) < d(a, f(a)) = g(a),
\]
contredisant la minimalité. Donc $f(a) = a$~; unicité comme
d'habitude (deux points fixes $a \neq b$ donnent $d(a,b) = d(f(a),
f(b)) < d(a,b)$).

Exemple non \omterm{def:b2:metric:compact}{compact}~: $f(x) = x + \frac1x$ sur $\intco{1}{+\infty}$~:
$\abs{f(x) - f(y)} = \abs{x - y}\,\abs{1 - \frac{1}{xy}} < \abs{x -
y}$ pour $x \neq y$ (car $xy > 1$), et pourtant $f(x) > x$ partout.
\end{solution}

\begin{solution}{exo:b2:metric:6}
Choisir $x_n \in K_n$. Tous les termes à partir du rang $n$ sont dans
$K_n$~; en particulier toute la suite est dans le \omterm{def:b2:metric:compact}{compact} $K_0$~:
extraire $x_{\varphi(k)} \to \ell$. Pour chaque $n$ fixé, les termes
$x_{\varphi(k)}$ avec $\varphi(k) \geq n$ sont dans le \emph{fermé}
$K_n$, donc la limite $\ell \in K_n$. Donc $\ell \in \bigcap K_n$.

Contre-exemple fermé~: $F_n = \intco{n}{+\infty}$ dans $\R$~:
emboîtés, fermés, non vides, intersection vide.
\end{solution}

\begin{solution}{exo:b2:metric:7}
La continuité de $f^{-1}$ signifie~: les images $f(F)$ des fermés $F
\subseteq K$ sont fermées (les images réciproques par $f^{-1}$ sont
les images par $f$). Un fermé $F$ dans le \omterm{def:b2:metric:compact}{compact} $K$ est \omterm{def:b2:metric:compact}{compact}
(\cref{thm:b2:metric:compactprops}~(1))~; son image \omterm{def:b2:metric:continuity}{continue} $f(F)$
est compacte, donc fermée. Donc $f^{-1}$ est \omterm{def:b2:metric:continuity}{continue}~: $f$ est un
homéomorphisme.

Contre-exemple~: $\gamma(t) = (\cos t, \sin t)$ de $\intco{0}{2\pi}$
(non \omterm{def:b2:metric:compact}{compact}) sur le cercle unité est une bijection \omterm{def:b2:metric:continuity}{continue}, mais
$\gamma^{-1}$ est discontinue en $(1,0)$~: les points du cercle juste
en dessous de l'axe ont des paramètres proches de $2\pi$, non de $0$.
\end{solution}

\begin{solution}{exo:b2:metric:8}
$C = \bigcap_n C_n$ où chaque $C_n$ (réunion de $2^n$ intervalles
fermés de longueur $3^{-n}$) est fermé~: $C$ est fermé et borné dans
$\R$, donc \omterm{def:b2:metric:compact}{compact} (\cref{thm:b2:metric:compactprops}~(2)).

Intérieur vide~: $C$ ne contient aucun intervalle de longueur $>
3^{-n}$ (il est dans $C_n$, dont les composantes ont cette longueur),
pour tout $n$.

Cardinalité~: les points de $C$ sont exactement les réels
$\sum_{n\geq1} a_n 3^{-n}$ avec chiffres $a_n \in \{0, 2\}$ (à chaque
étape, le tiers central supprimé retire le chiffre $1$)~;
l'application $(a_n) \mapsto \sum a_n 3^{-n}$ est une bijection de
$\{0,2\}^{\N^*}$ sur $C$ (injectivité comme dans le~%
\cref{exo:b2:structures:3}). Donc $C$ est équipotent à
$\{0,1\}^{\N}$~: non \omterm{def:b2:structures:countable}{dénombrable}, bien que de «~longueur nulle~».
\end{solution}

\begin{solution}{exo:b2:metric:9}
Supposons $a \notin f(X)$. L'image $f(X)$ est compacte (image
\omterm{def:b2:metric:continuity}{continue}), donc fermée~; ainsi
\[
\varepsilon = d\bigl(a, f(X)\bigr)
= \inf_{y \in f(X)} d(a, y) > 0
\]
(l'infimum d'une fonction \omterm{def:b2:metric:continuity}{continue} sur un \omterm{def:b2:metric:compact}{compact} est atteint~; s'il
valait $0$, $a$ serait adhérent au fermé $f(X)$, donc dedans).

Considérer l'orbite $x_n = f^n(a)$ ($x_0 = a$). Pour $p < q$~:
\[
d(x_p, x_q) = d\bigl(f^p(a), f^p(f^{q-p}(a))\bigr)
= d\bigl(a, f^{q-p}(a)\bigr) \geq \varepsilon,
\]
(isométrie itérée $p$ fois~; et $f^{q-p}(a) \in f(X)$ puisque $q - p
\geq 1$). Une suite à distances mutuelles $\geq \varepsilon$ n'a
aucune sous-suite convergente --- contredisant la compacité. Donc
$f(X) = X$.
\end{solution}

\begin{solution}{exo:b2:metric:10}
Soit $A, B \in GL_n(\C)$ et $p(z) = \det\bigl((1-z)A + zB\bigr)$~: un
polynôme en $z$ (chaque coefficient est affine en $z$~; le
\omterm{def:b2:linalg:det}{déterminant} est un polynôme en les coefficients). $p(0) = \det A \neq
0$~: $p$ n'est pas identiquement nul, donc il a un nombre fini de
racines $z_1, \dots, z_m$ (aucune égale à $0$ ni à $1$~: $p(1) = \det
B \neq 0$). Le plan $\C$ privé d'un nombre fini de points est \omterm{def:b2:metric:connected}{connexe}
par arcs~: un chemin de $0$ à $1$ évitant les $z_i$ existe (prendre
une ligne brisée passant par un point loin de toutes les racines, ou
un arc de cercle~; un nombre fini d'obstacles seulement). Le long
d'un tel chemin $\gamma$, $t \mapsto (1 - \gamma(t))A + \gamma(t)B$
est un chemin \omterm{def:b2:metric:continuity}{continu} \emph{à l'intérieur} de $GL_n(\C)$ de $A$ à $B$
(le \omterm{def:b2:linalg:det}{déterminant} ne s'y annule jamais). Donc $GL_n(\C)$ est \omterm{def:b2:metric:connected}{connexe}
par arcs --- contrairement à son cousin réel
(\cref{ex:b2:metric:glnr})~: le plan complexe a la place de
contourner les obstacles.
\end{solution}

\begin{solution}{exo:b2:metric:11}
\emph{Lipschitz~:} pour $a \in A$, $d(x, a) \leq d(x, y) + d(y,
a)$~; prendre l'infimum sur $a$~: $d(x, A) \leq d(x, y) + d(y, A)$,
et échanger $x, y$~: $\abs{d(x,A) - d(y,A)} \leq d(x,y)$.

\emph{Annulation~:} $d(x, A) = 0$ ssi il existe $a_n \in A$ avec
$d(x, a_n) \to 0$ ssi $x$ est limite de points de $A$ ssi $x \in
\overline A$.

\emph{L'interrupteur~:} pour des fermés disjoints $A, B$~: le
dénominateur $d(x,A) + d(x,B)$ ne s'annule jamais (cela forcerait $x
\in \overline A \cap \overline B = A \cap B = \emptyset$), donc
$\varphi$ est bien définie, et \omterm{def:b2:metric:continuity}{continue} comme quotient de fonctions
\omterm{def:b2:metric:continuity}{continues} de dénominateur ne s'annulant pas. $\varphi(x) = 0$ ssi
$d(x, A) = 0$ ssi $x \in A$~; $\varphi(x) = 1$ ssi $d(x, B) = 0$ ssi
$x \in B$~; et $0 \leq \varphi \leq 1$ partout.
\end{solution}

\begin{solution}{exo:b2:metric:12}
Soit $B(x_0, r_0)$ une boule quelconque~; nous y trouvons un point de
$\bigcap U_n$. Comme $U_1$ est dense et \omterm{def:b2:metric:topology}{ouvert}, $U_1 \cap B(x_0, r_0)$
est non vide et \omterm{def:b2:metric:topology}{ouvert}~: il contient une boule fermée $\overline
B(x_1, r_1)$ avec $0 < r_1 \leq \frac{r_0}2$ (rétrécir le rayon). Par
récurrence, $U_{n+1} \cap B(x_n, r_n)$ est non vide \omterm{def:b2:metric:topology}{ouvert}~: choisir
$\overline B(x_{n+1}, r_{n+1}) \subseteq U_{n+1} \cap B(x_n, r_n)$
avec $r_{n+1} \leq \frac{r_n}2$. Pour $p, q \geq n$, $x_p$ et $x_q$
sont dans $B(x_n, r_n)$ avec $r_n \leq 2^{-n}r_0$~: la suite est de
Cauchy, et converge vers un $\ell$ par complétude. Pour chaque $n$,
la queue de la suite est dans la boule \emph{fermée} $\overline
B(x_{n+1}, r_{n+1}) \subseteq U_{n+1} \cap B(x_0, r_0)$, donc $\ell
\in U_{n+1}$ pour tout $n$ et $\ell \in \overline B(x_1, r_1)
\subseteq B(x_0, r_0)$. Donc $\bigcap_n U_n$ rencontre toute boule~:
dense.

\emph{Application~:} si $\R = \bigcup_n F_n$ avec $F_n$ fermés
d'intérieur vide, alors $U_n = \R \setminus F_n$ sont denses
($\overline{U_n} = \R$ ssi $F_n$ est d'intérieur vide) et \omterm{def:b2:metric:topology}{ouverts}, et
Baire donne un point dans $\bigcap U_n = \R \setminus \bigcup F_n$~:
contradiction. En particulier $\R \neq \bigcup_{x \in D} \{x\}$ pour
un $D$ \omterm{def:b2:structures:countable}{dénombrable} (les singletons sont fermés d'intérieur vide)~:
$\R$ est non \omterm{def:b2:structures:countable}{dénombrable} --- \cref{thm:b2:structures:cantor} par une
autre voie.
\end{solution}

\begin{solution}{pb:b2:metric:1}
\textbf{1.} Soit $F$ fermé dans le \omterm{def:b2:metric:complete}{complet} $X$ et $(x_n) \subseteq F$
de Cauchy~: elle converge dans $X$ vers un $\ell$, et $\ell \in F$
car $F$ est fermé (les limites de suites de $F$ restent dans
$\overline F = F$)~: $F$ est \omterm{def:b2:metric:complete}{complet}. Réciproquement soit $A
\subseteq X$ \omterm{def:b2:metric:complete}{complet} et $x \in \overline A$~: une certaine suite de
$A$ converge vers $x$~; elle est de Cauchy, donc elle converge
\emph{dans} $A$~; les limites sont uniques, donc $x \in A$~: $A$
fermé. Comme $\bigl(C(I), d_\infty\bigr)$ est \omterm{def:b2:metric:complete}{complet}
(\cref{thm:b2:metric:rncomplete}), ses parties fermées sont
complètes.

\textbf{2.} Pour $y, z \in C(I)$ et $t \in I$~:
\[
\Bigl|\int_{t_0}^{t}y - \int_{t_0}^{t}z\Bigr| \leq
\abs{t - t_0}\,\sup_I\abs{y - z} \leq h\,d_\infty(y, z),
\]
et prendre le sup sur $t$.

\textbf{3.} En utilisant les deux équations de point fixe et
l'inégalité triangulaire~:
\[
d(\ell_g, \ell_{\widetilde g})
= d\bigl(g(\ell_g), \widetilde g(\ell_{\widetilde g})\bigr)
\leq d\bigl(g(\ell_g), g(\ell_{\widetilde g})\bigr) +
d\bigl(g(\ell_{\widetilde g}), \widetilde g(\ell_{\widetilde
g})\bigr)
\leq k\,d(\ell_g, \ell_{\widetilde g}) +
d\bigl(g(\ell_{\widetilde g}), \widetilde g(\ell_{\widetilde
g})\bigr),
\]
et résoudre en $d(\ell_g, \ell_{\widetilde g})$ (le coefficient $1 -
k$ est positif). La seconde inégalité majore l'écart évalué par
l'écart uniforme.

\textbf{4.} $g^m$ est une contraction sur un espace \omterm{def:b2:metric:complete}{complet} non
vide~: elle a un unique point fixe $\ell$
(\cref{thm:b2:metric:banach}). Alors $g^m(g(\ell)) = g(g^m(\ell)) =
g(\ell)$~: $g(\ell)$ est un point fixe de $g^m$, donc $g(\ell) =
\ell$ par unicité. Tout point fixe de $g$ en est un de $g^m$~:
unicité pour $g$. Orbites~: fixer $r \in \{0, \dots, m-1\}$~; la
sous-suite $(x_{qm + r})_q$ est la $g^m$-orbite partant de $x_r$,
donc elle converge vers $\ell$ quand $q \to \infty$ (Banach à
nouveau). Les $m$ sous-suites convergent vers le même $\ell$, donc
$x_n \to \ell$~: étant donné $\varepsilon$, chaque classe de résidus
est à terme à moins de $\varepsilon$, et il y a un nombre fini de
classes.

\textbf{5.} Si $y$ est \omterm{def:b2:metric:continuity}{continue} à valeurs dans $\intcc{y_0 - b}{y_0 +
b}$, l'intégrande $s \mapsto f(s, y(s))$ est \omterm{def:b2:metric:continuity}{continu} sur $I$
(composition), donc le membre de droite est $C^1$ de dérivée $f(t,
y(t))$ (théorème fondamental de l'analyse, volume de première année).
Si $y$ satisfait l'équation intégrale, elle est cette fonction $C^1$,
$y(t_0) = y_0$, et $y' = f(t, y)$. Réciproquement, en intégrant $y' =
f(s, y(s))$ de $t_0$ à $t$ on obtient l'équation intégrale.

\textbf{6.} $X_h$ contient la constante $y_0$~; il est fermé comme
image réciproque de $\intcc{0}{b}$ par l'application \omterm{def:b2:metric:continuity}{continue} $y
\mapsto d_\infty(y, y_0)$ (les distances sont $1$-lipschitziennes),
donc \omterm{def:b2:metric:complete}{complet} par la question 1. Stabilité~: pour $y \in X_h$ et $t
\in I$,
\[
\abs{T(y)(t) - y_0} = \Bigl|\int_{t_0}^{t}f(s, y(s))\dd s\Bigr|
\leq M\abs{t - t_0} \leq Mh \leq b ,
\]
la dernière étape par $h \leq b/M$ (ou $M = 0$, trivial). Et $T(y)$
est \omterm{def:b2:metric:continuity}{continue} ($C^1$ même, question 5)~: $T(y) \in X_h$.

\textbf{7.} Pour $t \in I$~:
\[
\abs{T(y)(t) - T(z)(t)} \leq \int_{t_0}^{t}\abs{f(s, y(s)) -
f(s, z(s))}\,\abs{\dd s} \leq L\abs{t - t_0}\,d_\infty(y,z)
\leq Lh\,d_\infty(y,z).
\]
Si $Lh < 1$~: $T$ est une contraction du \omterm{def:b2:metric:complete}{complet} non vide $X_h$, et
Banach donne un unique point fixe --- par la question 5, l'unique
solution.

\textbf{8.} Récurrence~; le cas $n = 1$ est l'inégalité médiane de la
question 7. En supposant la majoration pour $n$, pour $t \geq t_0$
(le cas $t \leq t_0$ est symétrique)~:
\[
\abs{T^{n+1}(y)(t) - T^{n+1}(z)(t)}
\leq L\int_{t_0}^{t}\abs{T^n(y)(s) - T^n(z)(s)}\dd s
\leq L\int_{t_0}^{t}\frac{L^n(s - t_0)^n}{n!}\dd s\;
d_\infty(y,z),
\]
et l'intégrale vaut $\frac{L^n(t - t_0)^{n+1}}{(n+1)!}$~: la
majoration avec $n + 1$. Donc $d_\infty(T^n y, T^n z) \leq
\frac{(Lh)^n}{n!}d_\infty(y, z)$, et $\frac{(Lh)^n}{n!} \to 0$ (la
série exponentielle converge)~: un certain $T^m$ est une contraction.
La question 4 s'applique sur le \omterm{def:b2:metric:complete}{complet} $X_h$~: $T$ a un unique point
fixe, c'est-à-dire le problème de Cauchy a exactement une solution
sur $I$ à valeurs dans $\intcc{y_0 - b}{y_0 + b}$.

\textbf{9.} Soit $y$ une solution du problème sur $I$ et supposons
l'ensemble $E = \{t \in I, t > t_0 : \abs{y(t) - y_0} > b\}$ non vide
(le côté $t < t_0$ est symétrique)~; soit $\tau = \inf E$. Par
continuité, $\abs{y(s) - y_0} \leq b$ pour $s \in \intcc{t_0}{\tau}$,
donc le graphe y est dans $R$, l'équation intégrale vaut jusqu'à
$\tau$, et
\[
\abs{y(\tau) - y_0} = \Bigl|\int_{t_0}^{\tau}f\bigl(s,
y(s)\bigr)\dd s\Bigr| \leq M(\tau - t_0) \leq Mh \leq b .
\]
Si $\tau < t_0 + h$, des points de $E$ arbitrairement proches de
$\tau$ par la droite donnent, par continuité, $\abs{y(\tau) - y_0}
\geq b$, donc $= b$~; mais alors l'affichage force $M(\tau - t_0) =
Mh$, c'est-à-dire $\tau = t_0 + h$~: contradiction. Donc $\tau = t_0
+ h$, $E \subseteq \{t_0 + h\}$, et l'affichage (en $\tau = t_0 + h$)
donne $\abs{y(t_0 + h) - y_0} \leq b$, contredisant l'appartenance à
$E$. Donc $E = \emptyset$~: toute solution sur $I$ reste dans la
bande, est un point fixe de $T$ dans $X_h$, et l'unicité est
inconditionnelle.

\textbf{10.} $T(y)(t) = 1 + \int_0^t y$. En partant de $y^{(0)}
\equiv 1$~:
\[
y^{(1)}(t) = 1 + t,\quad
y^{(2)}(t) = 1 + t + \frac{t^2}2,\quad\dots\quad
y^{(n)}(t) = \sum_{k=0}^{n}\frac{t^k}{k!}
\]
(récurrence~: intégrer la somme partielle ajoute le terme suivant).
Ce sont les sommes partielles de Taylor de $\eu^{\,t}$~; sur tout $I$
borné elles convergent uniformément vers $\eu^{\,t}$ (la queue est
dominée par la série numérique convergente $\sum h^k/k!$), qui est
bien l'unique solution.

\textbf{11.} $y \equiv 0$ est une solution. Pour $y_c$~: elle est
$C^1$ (les deux morceaux le sont, et en $t = c$ les dérivées
coïncident~: $0$ et $2(t - c) \to 0$), et pour $t > c$~: $y_c' = 2(t
- c) = 2\sqrt{(t-c)^2} = 2\sqrt{\abs{y_c}}$~; pour $t \leq c$ les
deux membres s'annulent. Donc le problème de Cauchy $y(0) = 0$ a une
infinité de solutions ($c \geq 0$ arbitraire, et $y \equiv 0$). Le
champ $\varphi(y) = 2\sqrt{\abs y}$ n'est pas \omterm{def:b2:metric:continuity}{lipschitzien} près de
$0$~: $\frac{\varphi(y) - \varphi(0)}{y - 0} = \frac{2}{\sqrt y} \to
+\infty$ quand $y \to 0^+$~: aucune constante $L$ ne convient sur un
voisinage de $0$ --- exactement là où toutes les solutions
bifurquent.

\textbf{12.} En séparant les variables (ou en vérifiant
directement), l'unique solution locale est $y(t) = \frac1{1 - t}$,
définie sur $\intoo{-\infty}{1}$ et explosant en $t = 1$. Pour le
rectangle $\intcc{-a}{a} \times \intcc{1 - b}{1 + b}$~: $M = \sup y^2
= (1 + b)^2$, donc la demi-largeur certifiée est $h =
\min\bigl(a, \frac{b}{(1+b)^2}\bigr)$. En maximisant
$\frac{b}{(1+b)^2}$~: dérivée nulle en $b = 1$, valeur $\frac14$.
Donc le théorème ne garantit la vie que sur
$\intcc{-\frac14}{\frac14}$ --- correctement moins que la vraie durée
de vie $1$ vers l'avant, et infiniment moins vers l'arrière~: le
théorème est local par nature, et l'explosion montre qu'il ne peut en
être autrement.

\textbf{13.} $g$ envoie $\intcc12$ dans lui-même~: $g$ décroît sur
$\intcc{1}{\sqrt2}$ et croît après (étudier $g'(x) = \frac12 -
\frac1{x^2}$), avec $g(1) = g(2) = \frac32$ et minimum $g(\sqrt2) =
\sqrt2 > 1$~: $g(\intcc12) \subseteq \intcc{\sqrt2}{\frac32}
\subseteq \intcc12$~; et $g$ envoie les rationnels sur des
rationnels. Contraction~: $\abs{g'(x)} = \abs{\frac12 - \frac1{x^2}}
\leq \frac12$ sur $\intcc12$ ($\frac1{x^2} \in
\intcc{\frac14}{1}$), donc l'inégalité des accroissements finis donne
$\abs{g(x) - g(y)} \leq \frac12\abs{x - y}$. Un point fixe vérifie
$\frac x2 = \frac1x$, c'est-à-dire $x^2 = 2$~: impossible dans $\Q$.
L'hypothèse en défaut est la complétude de $X$ ($\Q \cap \intcc12$
n'est pas \omterm{def:b2:metric:complete}{complet})~; dans le complété $\intcc12$ le point fixe est
$\sqrt2$ --- le théorème de Banach mené sur les rationnels
\emph{crée} l'irrationnel.

\textbf{14.} A posteriori~: $d(x_n, \ell) \leq d(x_n, x_{n+1}) +
d(x_{n+1}, \ell) \leq k\,d(x_{n-1}, x_n) + k\,d(x_n, \ell)$, d'où
$d(x_n, \ell) \leq \frac{k}{1-k}d(x_n, x_{n-1})$. Héron à partir de
$x_0 = \frac32$~: $x_1 = \frac{17}{12}$, $d(x_1, x_0) = \frac1{12}$,
$k = \frac12$~: la majoration a priori $\frac{k^n}{1-k}d(x_1,x_0) =
\frac{2^{-n+1}}{12}$ passe sous $10^{-6}$ pour la première fois en $n
= 18$. En réalité $x_1 = \frac{17}{12} \approx 1.41667$ (erreur
$2.5\cdot10^{-3}$), $x_2 = \frac{577}{408} \approx 1.4142157$ (erreur
$2.1\cdot10^{-6}$), $x_3 \approx 1.41421356237469$ (erreur
$1.6\cdot10^{-12}$)~: trois étapes suffisent. Chaque pas de Héron
\emph{élève au carré} l'erreur, grosso modo (convergence quadratique,
un phénomène de Newton~: \cref{ch:b2:realfun})~; l'estimation par
contraction, qui ne fait que la diviser par deux, est honnête pour le
pire cas mais pessimiste ici.

\textbf{15.} Appliquer la question 3 avec $g = T_{y_0}$ (une
$Lh$-contraction, $Lh < 1$) et $\widetilde g = T_{z_0}$, dont le
point fixe est $y[z_0]$. Pour tout $y$,
\[
\abs{T_{y_0}(y)(t) - T_{z_0}(y)(t)} = \abs{y_0 - z_0},
\]
(les intégrales sont identiques), donc $\sup_y
d_\infty(T_{y_0}(y), T_{z_0}(y)) = \abs{y_0 - z_0}$, et la question 3
donne
\[
d_\infty(y[y_0], y[z_0]) \leq \frac{\abs{y_0 - z_0}}{1 - Lh} .
\]

\textbf{16.} Sur chaque morceau, la question 15 appliquée avec les
valeurs à l'extrémité gauche comme données initiales majore l'écart à
l'extrémité droite~:
\[
d_\infty \leq \frac{1}{1 - 1/2}\,(\text{écart à gauche})
= 2\,(\text{écart à gauche}) .
\] Par
récurrence sur les $m$ morceaux, l'écart final est au plus
$2^m\abs{y_0 - z_0}$, et l'écart uniforme sur tout le segment obéit à
la même majoration (le sup de chaque morceau est contrôlé à son
stade). Avec des morceaux de longueur $h \asymp \frac1{2L}$, le
facteur est $2^m = 2^{\,\text{longueur}\cdot 2L}$~: exponentiel en la
longueur de l'intervalle, exactement comme le prédit la majoration de
Gronwall $\eu^{L\abs{t-t_0}}$, avec de meilleures constantes.

\textbf{17.} Même schéma~: pour $y \in X_h$,
\[
\abs{T_f(y)(t) - T_g(y)(t)} \leq \int_{t_0}^t \abs{f(s,y(s)) -
g(s,y(s))}\,\abs{\dd s} \leq \varepsilon h ,
\]
donc la question 3 (avec $T_f$ la contraction, $T_g$ l'application
perturbée) donne $d_\infty \leq \frac{\varepsilon h}{1 - Lh}$.

\textbf{18.} Par la question 17 appliquée à $f = f_\lambda$, $g =
f_\mu$~: $d_\infty(y_\lambda, y_\mu) \leq \frac{Ch}{1 -
Lh}\abs{\lambda - \mu}$~: l'application solution est \omterm{def:b2:metric:continuity}{lipschitzienne}
de constante $\frac{Ch}{1-Lh}$.

\textbf{19.} Chaque argument n'a utilisé que~: les axiomes
métriques, la complétude de l'étage, la majoration $\abs{\int} \leq
\int\abs{\cdot}$, et la propriété de Lipschitz de $f$ --- toutes
disponibles pour les fonctions à valeurs dans $\R^n$ avec $d_\infty$
bâtie sur l'une des distances équivalentes du~%
\cref{ex:b2:metric:examples} (complètes par
\cref{thm:b2:metric:rncomplete}). Pour $y' = Ay$ avec $A$
nilpotente~: $y^{(0)} \equiv (c_1, c_2)$,
\[
y^{(1)}(t) = (c_1, c_2) + \int_0^t (c_2, 0)\,\dd s
= (c_1 + tc_2,\; c_2),
\]
et $Ay^{(1)}(s) = (c_2, 0)$ à nouveau~: $y^{(2)} = y^{(1)}$.
L'itération est stationnaire à partir de $n = 1$, à la solution
exacte $y(t) = (c_1 + tc_2, c_2) = \eu^{tA}y(0)$ --- la nilpotence
tronque la série exponentielle, et Picard le remarque.

\textbf{20.} Fixer $y \in X$ et poser $g_y(x) = y - \eta(x)$~: une
$k$-contraction du \omterm{def:b2:metric:complete}{complet} $X$, donc il existe exactement un $x$ avec
$x + \eta(x) = y$~: l'application $\Phi = \mathrm{id} + \eta$ est
bijective. Réciproque \omterm{def:b2:metric:continuity}{lipschitzienne}~: si $\Phi(x) = y$ et $\Phi(x')
= y'$,
\[
d(x, x') \leq d(y, y') + d\bigl(\eta(x), \eta(x')\bigr)
\leq d(y, y') + k\,d(x, x'),
\]
donc $d(x, x') \leq \frac{1}{1-k}d(y, y')$. (Les distances ici
proviennent de la structure de norme, donc $d(a - c, b - c) = d(a,
b)$, ce qu'a utilisé la première inégalité.)

\textbf{21.} $g(x) = m + e\sin x$ est $e$-lipschitzienne sur le
\omterm{def:b2:metric:complete}{complet} $\R$ (inégalité des accroissements finis, $\abs{g'} =
\abs{e\cos x} \leq e < 1$)~: Banach donne une unique solution et la
convergence globale de l'itération. Pour $e = \frac12$, $m = 1$, $x_0
= 1$~: $x_1 = 1 + \frac12\sin 1 \approx 1.42074$, $d(x_1, x_0)
\approx 0.4207$, et la majoration a priori $\frac{(1/2)^n}{1/2}\cdot
0.4207 \leq 10^{-3}$ vaut pour la première fois en $n = 10$~: dix
itérations certifiées (la vraie valeur $x \approx 1.4987$ est en fait
atteinte à $10^{-3}$ près déjà vers $n = 5$).

\textbf{22.} Utiliser la description par chiffres
(\cref{exo:b2:metric:8})~: $C$ est l'ensemble des sommes
$\sum_{n\geq1}a_n3^{-n}$, $a_n \in \{0, 2\}$. Alors $S_1(C) = \{x/3 :
x \in C\}$ est le sous-ensemble avec $a_1 = 0$, et $S_2(C) = \{x/3 +
2/3\}$ le sous-ensemble avec $a_1 = 2$~: leur réunion, sur le choix
libre de $a_1$, est exactement $C$. Unicité en une phrase~: sur
l'espace des parties \omterm{def:b2:metric:compact}{compactes} non vides de $\intcc01$ métrisé par la
distance de Hausdorff, $A \mapsto S_1(A) \cup S_2(A)$ est une
$\frac13$-contraction d'un espace \omterm{def:b2:metric:complete}{complet}, donc Banach n'autorise
qu'un seul ensemble fixe --- le volume de troisième année rend la
métrique de Hausdorff et cet argument rigoureux.

\textbf{23.} Soit $Z = \{t \in J : y(t) = z(t)\}$~: non vide ($t_0
\in Z$), fermé dans $J$ (égaliseur de deux applications \omterm{def:b2:metric:continuity}{continues}~:
image réciproque de $\{0\}$ par $y - z$). \omterm{def:b2:metric:topology}{Ouvert}~: si $t_1 \in Z$,
appliquer le théorème local (question 8) au point $(t_1, y(t_1))$,
dans un rectangle où $f$ est \omterm{def:b2:metric:continuity}{lipschitzienne}~: sur un petit intervalle
autour de $t_1$, $y$ et $z$ résolvent le même problème de Cauchy,
donc coïncident là (l'unicité inconditionnelle de la question 9)~: un
voisinage de $t_1$ est dans $Z$. Une partie non vide de l'intervalle
$J$ à la fois ouverte et fermée dans $J$ est $J$ tout entier (les
intervalles sont \omterm{def:b2:metric:connected}{connexes}, \cref{thm:b2:metric:connectedness})~: $y =
z$ sur $J$.

\textbf{24.} Ici $f(t, y) = \alpha(t)y + \beta(t)$ est
$L$-lipschitzienne en $y$ sur tout $\intcc AB \times \R$ avec $L =
\sup\abs\alpha$ (finie~: $\alpha$ \omterm{def:b2:metric:continuity}{continue} sur un segment), et
aucune bande $\intcc{y_0 - b}{y_0 + b}$ n'est nécessaire~: prendre $X
= C(\intcc AB)$ entier, sur lequel $T$ est bien défini. La récurrence
de la question 8 tourne mot pour mot et donne $d_\infty(T^ny, T^nz)
\leq \frac{(L(B - A))^n}{n!}d_\infty(y, z)$~: un itéré est une
contraction quelle que soit la longueur $B - A$, et la question 4
conclut~: une et une seule solution sur le segment entier. La
linéarité entre exactement une fois~: elle rend la majoration de
Lipschitz globale en $y$, supprimant le rectangle et son $h \leq
b/M$.

\textbf{25.} La \emph{complétude} a transformé la suite de Cauchy des
itérés en une solution effective (questions 1, 6, 8), et son absence
a laissé $\sqrt2$ s'échapper de $\Q$ (question 13). La
\emph{contraction} a donné l'unicité, l'algorithme, et les barres
d'erreur (questions 7, 14). \emph{L'astuce de l'itéré} a supprimé la
condition de petitesse $Lh < 1$, de sorte que l'intervalle certifié
ne dépend que de $M$, non de $L$ --- et a rendu les équations
linéaires globales (questions 8, 24). La \emph{connexité} a promu
l'unicité locale en unicité globale (question 23). Les
contre-exemples~: $2\sqrt{\abs y}$ garde Lipschitz (question 11),
$y^2$ garde la localité (question 12), $\Q$ garde la complétude
(question 13). Le sommet est le théorème de Picard--Lindelöf
(question 8)~; lorsque $f$ est seulement \omterm{def:b2:metric:continuity}{continue}, l'existence
survit mais l'unicité meurt, et la démonstration échange la
contraction contre la compacité d'ensembles de fonctions --- le
théorème de Peano via Arzelà--Ascoli, dans le volume de troisième
année.
\end{solution}
